\chapter{Verification}
\label{ch:pebble-verification}

An interpreter that follows an invalid operand ID can fail before it has a chance to explain the malformed program.
A verifier must therefore check the safety of its own queries in a deliberate order.
PebbleIR first establishes bounds and graph shape, then record shape, then operand availability and types.

\section{Layers of trust}
\label{sec:pebble-verification-layers}
The verifier accepts a valid initialized C object, not an arbitrary byte buffer cast to \code{Pir}.
Serialized input enters through the parser.
It rejects zero blocks, excessive counts, unknown terminators, invalid targets, edges into entry, and duplicate targets of a conditional branch.
It builds predecessor sets and requires every declared block to be reachable from entry.
This last rule keeps dead regions from complicating dominance and is stricter than many production IRs.

Next it checks every value header before following any value operand.
That ordering matters when a valid-looking operand ID points to a record whose block field is out of bounds.
A dominance lookup using that block before checking it would make the verifier itself unsafe.
The malformed fixtures include exactly this boundary.

\begin{algorithm}[htbp]
\caption{Verifying in dependency order.}
\label{alg:pebble-verify}
\begin{pseudocode}
\Require initialized function object $F$
\Ensure verified $F$ and fresh facts, or diagnostic with facts unchanged
\State validate counts and all terminators before following edges
\State derive predecessors; compute reachability; reject unreachable blocks
\State compute dominators over the reachable graph
\State validate every value's opcode, type, block, and canonical fields
\State validate phi prefixes and exact predecessor coverage
\State check live operands, use-point dominance, and operand types
\State check return and branch operands at block terminators
\State publish facts tagged with owner and revision
\end{pseudocode}
\end{algorithm}

Algorithm~\ref{alg:pebble-verify} is implemented by \code{pir\_verify}.
Facts are accumulated locally and copied to the output only on success.
This keeps a caller from mistaking a partly computed dominator table for a usable analysis result after a rejected function.

\section{Uses belong to program points}
\label{sec:pebble-use-points}
An ordinary operand in block $B$ must be defined earlier in $B$, or in a different block that dominates $B$.
A terminator operand can use any earlier live value in its block.
A phi operand is checked at the corresponding predecessor's terminator, not at the phi's textual location.
Listing~\ref{lst:pebble-availability} is the common availability query; its caller supplies the appropriate block and ordering boundary.

\begin{listing}[htbp]
\inputminted[firstline=50,lastline=56]{c}{example-code/pebbleir/pebbleir.c}
\caption{Availability combines live-ID validation, block dominance, and same-block order.}
\label{lst:pebble-availability}
\end{listing}

Suppose a diamond defines value 5 in its left branch and value 7 in its right branch.
Changing an operand inside the left branch to value 7 preserves both the numeric ID range and the byte type.
It nevertheless fails availability because the right block does not dominate the left block.
A phi in the join may legitimately use both values, each on its own incoming edge.
This is why a type checker alone is not an SSA verifier.

\section{What acceptance does not establish}
\label{sec:pebble-verify-limits}
Verification establishes internal well-formedness, not equivalence to a source program.
A function that returns constant zero is valid even when the source expression should return 128.
It also does not establish termination, profitability, or absence of implementation defects in the verifier itself.
Differential tests and local correctness arguments cover different obligations.

Let $N$ be allocated values, $B$ blocks, and $R$ the number of dominator sweeps.
With the implementation's one-word block sets, the dominator portion costs $O(RB^2)$ and the later scans cost $O(N+B)$.
The interface recomputes facts rather than trusting stale caches; frequent verification therefore has a visible compile-time cost.
Production builds can choose checkpoint frequency without changing which invariants clients are allowed to assume.

\section{Exercises}
\label{sec:pebble-eight-exercises}
\begin{enumerate}
\item \label{ex:pebble-welltyped-bad} Construct the sibling-branch misuse described above and identify the check that rejects it.
\item \label{ex:pebble-phi-use} Explain why a loop-header phi can refer to a latch value with a larger ID, while an ordinary header addition cannot use that latch value directly.
\end{enumerate}

\section{Further reading}
\label{sec:pebble-eight-reading}
Cytron and colleagues supply the dominance foundations behind SSA's definition-use discipline~\cite{cytron-1991-ssa}.
Safonov provides context for treating representation checks as part of a compiler's trust argument rather than as optional debugging decoration~\cite{safonov-2010-trustworthy}.
